\leanverified{paper\_conclusion\_screen\_minimal\_exact\_extension\_contraction\_basis}
\begin{theorem}[部分屏幕最小精确上扩的收缩拟阵分类]\label{thm:conclusion-screen-minimal-exact-extension-contraction-basis}
沿用定理 \ref{thm:conclusion-screen-visible-external-audit-conservation} 的记号。对任意部分屏幕
\[
S\subseteq \overrightarrow{\mathcal F_m^{n-1}},
\]
定义其最小精确上扩族
\[
\mathcal E_{\min}(S)
:=
\Bigl\{
T\subseteq \overrightarrow{\mathcal F_m^{n-1}}:\ 
S\subseteq T,\ a(T)=0,\ |T\setminus S|=a(S)
\Bigr\}.
\]
则存在自然双射
\[
\mathcal E_{\min}(S)\cong \mathcal B(M_{m,n}/S).
\]
更精确地，\(T\in \mathcal E_{\min}(S)\) 当且仅当
\[
T=S\sqcup B,
\qquad
B\text{ 是 }M_{m,n}/S\text{ 的一组基}.
\]
\end{theorem}
\begin{proof}
由定理 \ref{thm:conclusion-screen-visible-external-audit-conservation}，
\[
a(S)=\rho_{m,n}-r(S)=r(M_{m,n}/S).
\]
若 \(T\supseteq S\)，记
\[
B:=T\setminus S\subseteq \overrightarrow{\mathcal F_m^{n-1}}\setminus S.
\]
则
\[
a(T)=0
\iff
r(T)=\rho_{m,n}
\iff
r_{M_{m,n}/S}(B)=r(M_{m,n}/S).
\]
因此，\(T\) 为精确上扩当且仅当 \(B\) 在收缩拟阵 \(M_{m,n}/S\) 中张成满秩。

再要求
\[
|T\setminus S|=a(S)=r(M_{m,n}/S),
\]
便等价于 \(B\) 既张成又具有最小可能基数，也即 \(B\) 是 \(M_{m,n}/S\) 的一组基。反向同理。故 \(T\mapsto T\setminus S\) 给出所述自然双射。证毕。
\end{proof}

\begin{theorem}[精确化缺口配分函数的 Tutte 特化]\label{thm:conclusion-screen-deficit-partition-function-tutte-specialization}
沿用定理 \ref{thm:conclusion-screen-minimal-exact-extension-contraction-basis} 的记号。定义部分屏幕 \(S\) 的精确化配分函数
\[
Z_S(z):=
\sum_{U\subseteq \overrightarrow{\mathcal F_m^{n-1}}\setminus S}
z^{\,a(S\cup U)}.
\]
则有精确闭式
\[
Z_S(z)=T_{M_{m,n}/S}(z+1,2),
\]
其中 \(T_{M_{m,n}/S}\) 为收缩拟阵 \(M_{m,n}/S\) 的 Tutte 多项式。特别地，
\[
|\mathcal E_{\min}(S)|=T_{M_{m,n}/S}(1,1).
\]
\end{theorem}
\begin{proof}
对任意
\[
U\subseteq \overrightarrow{\mathcal F_m^{n-1}}\setminus S,
\]
由收缩拟阵的秩公式与定理 \ref{thm:conclusion-screen-visible-external-audit-conservation}，
\[
a(S\cup U)
=
\rho_{m,n}-r(S\cup U)
=
r(M_{m,n}/S)-r_{M_{m,n}/S}(U).
\]
故
\[
Z_S(z)
=
\sum_{U\subseteq E(M_{m,n}/S)}
z^{\,r(M_{m,n}/S)-r_{M_{m,n}/S}(U)}.
\]
而 Tutte 多项式定义为
\[
T_M(x,y)
=
\sum_{A\subseteq E(M)}
(x-1)^{r(M)-r(A)}(y-1)^{|A|-r(A)}.
\]
令 \(y=2\) 时，第二因子恒为 \(1\)，遂得
\[
T_M(z+1,2)=\sum_{A\subseteq E(M)} z^{\,r(M)-r(A)}.
\]
取 \(M=M_{m,n}/S\) 即得所示公式。

最后，收缩拟阵的基数目满足标准恒等式
\[
|\mathcal B(M_{m,n}/S)|=T_{M_{m,n}/S}(1,1).
\]
再由定理 \ref{thm:conclusion-screen-minimal-exact-extension-contraction-basis}，
\[
|\mathcal E_{\min}(S)|=T_{M_{m,n}/S}(1,1).
\]
证毕。
\end{proof}

\leanverified{paper\_conclusion\_screen\_deficit\_unit\_descent\_chain}
\begin{theorem}[精确化缺口的一步一单位下降链]\label{thm:conclusion-screen-deficit-unit-descent-chain}
对任意部分屏幕
\[
S\subseteq \overrightarrow{\mathcal F_m^{n-1}},
\]
存在一条单调上升链
\[
S=S_0\subset S_1\subset \cdots \subset S_{a(S)}
\]
满足
\[
|S_j\setminus S|=j,
\qquad
a(S_j)=a(S)-j
\qquad
(0\le j\le a(S)),
\]
且终点 \(S_{a(S)}\) 为精确屏幕。
\end{theorem}
\begin{proof}
由定理 \ref{thm:conclusion-screen-minimal-exact-extension-contraction-basis}，可取
\[
B=\{b_1,\dots,b_{a(S)}\}
\]
为收缩拟阵 \(M_{m,n}/S\) 的一组基。定义
\[
S_j:=S\cup\{b_1,\dots,b_j\}.
\]
因为基的任意前缀在收缩拟阵中都独立，故
\[
r_{M_{m,n}/S}(\{b_1,\dots,b_j\})=j.
\]
于是
\[
r(S_j)=r(S)+j,
\]
从而
\[
a(S_j)=\rho_{m,n}-r(S_j)=\rho_{m,n}-r(S)-j=a(S)-j.
\]
当 \(j=a(S)\) 时，
\[
r(S_{a(S)})=r(S)+a(S)=\rho_{m,n},
\]
故 \(S_{a(S)}\) 为精确屏幕。证毕。
\end{proof}

\begin{corollary}[单洞缺口的一步闭合判据]\label{cor:conclusion-screen-single-hole-one-step-closure}
对任意部分屏幕 \(S\subseteq \overrightarrow{\mathcal F_m^{n-1}}\)，以下命题等价：
\begin{enumerate}
  \item \(a(S)=1\)；
  \item \(\beta_{\mathrm{top}}(S)=1\)；
  \item 存在某张面 \(\vec f\notin S\) 使
        \[
        a(S\cup\{\vec f\})=0.
        \]
\end{enumerate}
\end{corollary}
\begin{proof}
由定理 \ref{thm:conclusion-screen-exactification-triple-identification}，
\[
a(S)=\beta_{\mathrm{top}}(S),
\]
故前两条等价。

若 \(a(S)=1\)，则定理 \ref{thm:conclusion-screen-deficit-unit-descent-chain} 给出一条长度为 \(1\) 的下降链，故存在某张面 \(\vec f\notin S\) 使
\[
a(S\cup\{\vec f\})=0.
\]
反之，若存在 \(\vec f\notin S\) 满足 \(a(S\cup\{\vec f\})=0\)，则由定理 \ref{thm:conclusion-screen-rank-gain-closure-diminishing-returns} 的单面 \(0\text{--}1\) 律，
\[
a(S)-a(S\cup\{\vec f\})\in\{0,1\}.
\]
既然 \(a(S\cup\{\vec f\})=0\) 且 \(S\) 尚未精确，便只能有 \(a(S)=1\)。证毕。
\end{proof}

\leanverified{paper\_conclusion\_screen\_basis\_exchange\_graph\_geodesic\_rigidity}
\begin{theorem}[极小精确屏幕交换图的测地刚性]\label{thm:conclusion-screen-basis-exchange-graph-geodesic-rigidity}
记全部极小精确屏幕之集合为
\[
\mathcal B_{m,n}:=\mathcal B(M_{m,n}).
\]
定义其 basis-exchange 图 \(\Gamma_{\mathrm{ex}}\)：若
\[
B'=(B\setminus\{e\})\cup\{f\}
\]
对某个 \(e\in B\)、\(f\notin B\) 成立，则连边 \(B\sim B'\)。则对任意
\[
B,B'\in \mathcal B_{m,n},
\]
有精确距离公式
\[
d_{\Gamma_{\mathrm{ex}}}(B,B')
=
\frac{|B\triangle B'|}{2}
=
\frac{d_m^{\mathrm{arith}}(\mathcal G(B),\mathcal G(B'))}{2}.
\]
\end{theorem}
\begin{proof}
一次基交换恰移去一个元素并加入一个元素，故对称差大小至多减少 \(2\)。因此任意从 \(B\) 到 \(B'\) 的交换链长度都至少为
\[
\frac{|B\triangle B'|}{2}.
\]

反过来，若 \(B\neq B'\)，取任意
\[
e\in B\setminus B'.
\]
由强基交换定理，存在
\[
f\in B'\setminus B
\]
使得
\[
B_1:=(B\setminus\{e\})\cup\{f\}
\]
仍为基。并且
\[
|B_1\triangle B'|=|B\triangle B'|-2.
\]
对 \(|B\triangle B'|\) 归纳，便构造出一条长度恰为 \(|B\triangle B'|/2\) 的交换链。

最后，由定理 \ref{thm:conclusion-screen-audit-gap-arithmetic-lipschitz} 中
\[
d_m^{\mathrm{arith}}(\mathcal G(B),\mathcal G(B'))=|B\triangle B'|
\]
即得第二个等号。证毕。
\end{proof}

\begin{theorem}[朝向 exact sector 的有向算术距离公式]\label{thm:conclusion-screen-directed-arithmetic-distance-to-exact}
对任意部分屏幕
\[
S\subseteq \overrightarrow{\mathcal F_m^{n-1}},
\]
定义
\[
\mathcal E(S):=
\Bigl\{
T\subseteq \overrightarrow{\mathcal F_m^{n-1}}:\ 
S\subseteq T,\ a(T)=0
\Bigr\}.
\]
则有精确公式
\[
a(S)=\min_{T\in\mathcal E(S)} d_m^{\mathrm{arith}}(\mathcal G(S),\mathcal G(T)).
\]
并且等号成立当且仅当 \(T\) 是 \(S\) 的最小精确上扩。
\end{theorem}
\leanverified{paper\_conclusion\_screen\_directed\_arithmetic\_distance\_to\_exact}
\begin{proof}
若 \(T\in\mathcal E(S)\)，则 \(T\supseteq S\)，故
\[
d_m^{\mathrm{arith}}(\mathcal G(S),\mathcal G(T))
=
|T\triangle S|
=
|T\setminus S|.
\]
由 \(a(T)=0\) 即 \(r(T)=\rho_{m,n}\)，可知 \(T\setminus S\) 在收缩拟阵 \(M_{m,n}/S\) 中张成满秩，因此其基数必满足
\[
|T\setminus S|\ge r(M_{m,n}/S)=a(S).
\]
于是
\[
d_m^{\mathrm{arith}}(\mathcal G(S),\mathcal G(T))\ge a(S)
\]
对所有 \(T\in\mathcal E(S)\) 成立。

另一方面，定理 \ref{thm:conclusion-screen-minimal-exact-extension-contraction-basis} 给出某个最小精确上扩
\[
T_\ast\supseteq S
\]
满足
\[
|T_\ast\setminus S|=a(S),
\]
故
\[
d_m^{\mathrm{arith}}(\mathcal G(S),\mathcal G(T_\ast))=a(S).
\]
于是最小值公式成立，而取等类正是所有最小精确上扩。证毕。
\end{proof}

\begin{corollary}[极小精确屏幕算术码的常权偶距结构]\label{cor:conclusion-screen-minimal-exact-code-constant-weight-even-metric}
定义极小精确屏幕的算术码
\[
\mathcal C_{m,n}
:=
\{\mathcal G(B):\ B\in \mathcal B_{m,n}\}.
\]
则：
\begin{enumerate}
  \item \(\mathcal C_{m,n}\) 是常权码，所有码字支撑大小都等于
        \[
        \rho_{m,n}=2^{mn};
        \]
  \item 任意两码字之间的 arithmetic 距离都为偶数；
  \item 若 \(|\mathcal B_{m,n}|>1\)，则其最小非零距离恰为 \(2\)。
\end{enumerate}
\end{corollary}
\leanverified{paper\_conclusion\_screen\_minimal\_exact\_code\_constant\_weight\_even\_metric}
\begin{proof}
由定理 \ref{thm:conclusion-screen-zero-slack-matroid-basis-criterion}，极小精确屏幕恰为边界行拟阵 \(M_{m,n}\) 的全部基，而每个基的基数都等于
\[
\rho_{m,n}=2^{mn}.
\]
故 \(\mathcal C_{m,n}\) 为常权码。

若 \(B,B'\) 为两组基，则
\[
|B|=|B'|=\rho_{m,n},
\]
从而
\[
|B\triangle B'|
=
|B\setminus B'|+|B'\setminus B|
=
2|B\setminus B'|
\]
必为偶数。再由 arithmetic 距离与对称差大小一致，第二条成立。

若 \(|\mathcal B_{m,n}|>1\)，取任意不同基 \(B\neq B'\)。由强基交换定理，存在
\[
e\in B\setminus B',
\qquad
f\in B'\setminus B
\]
使得
\[
(B\setminus\{e\})\cup\{f\}
\]
仍为一组基，且与 \(B\) 的对称差恰为 \(2\)。故最小非零距离至多为 \(2\)。另一方面，常权码不可能出现距离 \(1\)，故最小非零距离恰为 \(2\)。证毕。
\end{proof}

\endinput
